% Section 5
\section{Signs, reinforcement and monotonicity}\label{sec:signs}

Monotone ecological inference~\citep{egghm2025} replaces the functional form of
LCE by beliefs about the direction of two associations: whether a group is more
likely than others to choose an outcome within the same unit (the \emph{between}
association), and whether its rate rises with its share across units (the
\emph{within}, or contextual, association). This section states those beliefs in
R$\times$C, shows how they bound the target on their own and on top of a linear
rung, and adds a nonparametric monotone rung.

Two settings are used. In the \emph{finite-population} setting the design is
fixed, unit weights are $p_i=N_i/\sum_jN_j$, $\E_p[f]=\sum_ip_if_i$,
$\mu_r=\E_p[X_r]>0$, and the realised rates satisfy $\beta_i\in S(X_i,T_i)$. The
target is the actual rate $B_{rc}=\sum_iw_{ir}\beta_{irc}$, with
$w_{ir}=p_iX_{ir}/\mu_r$ as before, and the rates are read as real numbers (cell
masses need not be whole numbers of people). In the \emph{superpopulation}
setting, $m(v)=\E[\beta\mid V=v]$, the target is $\Bt$, and a superscript $^*$ on
an association means it is computed with $m$ in place of $\beta$. In this section
$D$ without a unit index denotes a gap between two groups' rates, not a district.

\subsection{The associations and their identities}

Fix a column and drop it from the notation. Write $\Cov_w$ for a covariance under
weights $w$, $w^r_i=w_{ir}$, and $w^{\neg r}_i=p_i(1-X_{ir})/(1-\mu_r)$.

\begin{definition}[Associations]\label{def:assoc}
\begin{enumerate}[label=(\roman*)]
\item Between, pairwise: $\delta_B(r,s)=\E_p[X_rX_s(\beta_r-\beta_s)]$.
\item Between, one-vs-rest: $\delta^r_B=\sum_{s\ne r}\delta_B(r,s)$.
\item Within, pooled (W2): $\Delta=\sum_s\mu_s\Cov_{w^s}(X,\beta_s)\in\R^R$.
\item Within, one-vs-rest (W1): with
$\beta^{\neg r}_i=\sum_{s\ne r}X_{is}\beta_{is}/(1-X_{ir})$, the rate of everyone
outside row $r$,
$\delta^r_W=\mu_r\Cov_{w^r}(\beta_r,X_r)+(1-\mu_r)\Cov_{w^{\neg r}}(\beta^{\neg r},X_r)$.
\end{enumerate}
\end{definition}
\begin{reading}
The between association says that, on average within the same unit, group $r$
is more likely than the others to choose the outcome. It is a weighted average
with weights $p_iX_{ir}(1-X_{ir})$: pure units get no weight and a large mixed
unit can dominate. It is not an ordering in every unit, nor against each rival.
The within association says that rates rise, on average, with group $r$'s share.
In W1 the predictor in both covariances is $X_r$, and the complement's rate is
that of a changing mix of the other groups, so either covariance can be negative
when their sum is positive. A sign on $\Delta_r$ signs one pooled sum of
group-specific covariances, not each group. W1 and W2 can have opposite signs
while every group's rate rises: with groups A, B, C, two units of equal size with
shares $(.2,.7,.1)$ and $(.6,.05,.35)$, and rates $(.40,.80,.10)$ and
$(.42,.82,.12)$, every rate rises by two points, yet $\delta^A_W=-0.02633$ and
$\Delta_A=+0.00110$. Covariances also reflect sorting; a causal reading needs a
separate design. W1 is the one-vs-rest collapse suggested by EGGHM; W2, the
pairwise weighting and reinforcement with those pairings are extensions proved
here.
\end{reading}

Four data-computable quantities appear: Goodman's no-intercept regression
$g=M^{-1}\E_p[XT]$ with $M=\E_p[XX^\top]$; the neighbourhood model
$\NM_r=\E_p[X_rT]/\mu_r$; the one-vs-rest regression slope
$\ER^r=\Cov_p(T,X_r)/\Var_p(X_r)$; and $\NM_{\neg r}=\E_p[(1-X_r)T]/(1-\mu_r)$.
Write $D^r=B_r-B_{\neg r}$ for the gap between group $r$ and the rest, where
$B_{\neg r}=\sum_{s\ne r}\mu_sB_s/(1-\mu_r)$.

\begin{lemma}[Accounting identities]\label{lem:M1}\status{proved}
Assume A1. (a) $\one^\top\Delta=0$ and, if $M$ is invertible, $g-B=M^{-1}\Delta$.
(b) $\NM_r-B_r=-\delta^r_B/\mu_r$. (c) If $\Var_p(X_r)>0$ and $0<\mu_r<1$, the
collapsed population $(X_r,\beta_r,\beta^{\neg r})$ is a valid two-group
population, so
\[
\ER^r-D^r=\frac{\delta^r_W}{\Var_p(X_r)},\qquad(\NM_r-\NM_{\neg r})-D^r=-\frac{\delta^r_B}{\mu_r(1-\mu_r)}.
\]
(d) $B$, $\delta_B(r,s)$, $\delta^r_B$, $\Delta$ and $\delta^r_W$ are affine in
the unit rates, with coefficients that depend only on the design.
\end{lemma}
\begin{reading}
These are accounting identities. The neighbourhood model misses the true gap
exactly by the between association, and the aggregate regression slope misses it
exactly by the within association. So a sign belief about either association is
a one-sided bound relative to a number computed from the data. Signing $\Delta_r$
does not in general sign the bias of row $r$'s Goodman coefficient, because of
the matrix $M^{-1}$.
\end{reading}

\subsection{Sign restrictions on realised tables (N1)}

\begin{definition}[N1 feasible set]
A \emph{sign restriction} is $g_k(\beta)\ge-\varepsilon_k$, where $g_k$ is plus or
minus one of the associations of Definition~\ref{def:assoc} (for any column) and
$\varepsilon_k\ge0$ is a slack (default 0). Optional unit-level restrictions
$\beta_i\in U_i$, with $U_i$ polyhedral, express unit-level orderings or
structural zeros. The \emph{N1 feasible set} is
$F=\{\beta:\beta_i\in S(X_i,T_i)\cap U_i\ \forall i;\ g_k(\beta)\ge-\varepsilon_k\ \forall k\}$.
\emph{Contextual reinforcement} (CR) for a chosen pair $(g_B,g_W)$ of a between
and a within association is $F_{CR}=F_+\cup F_-$ with
$F_\pm=F\cap\{\pm g_B\ge0,\ \pm g_W\ge0\}$: the two associations have the same
sign, whichever it is.
\end{definition}

\begin{proposition}[N1 is sharp for continuous unit tables]\label{prop:N1}\status{proved}
Let $\varphi=\sum a_{rc}B_{rc}$ and read the unit tables as non-negative real cell
masses. (a) The set of values of $\varphi$ attained by continuous unit tables that
reproduce the margins and satisfy the restrictions is $\varphi(F)$; if
$F\ne\emptyset$ it equals $[\min_F\varphi,\max_F\varphi]$. (b) $F=\emptyset$ iff
the restrictions contradict the margins in this continuous problem. (c) Under CR
the set is $\varphi(F_+)\cup\varphi(F_-)$, a union of at most two intervals.
(d) $\varepsilon\mapsto\varphi(F(\varepsilon))$ is nested componentwise, with
continuous piecewise-linear endpoints where $F(\varepsilon)\neq\emptyset$. With
several slacks, the set of slack vectors that preserves a given conclusion is in
general not a product of intervals; its boundary is the breakdown frontier. A
scalar breakdown value needs a stated conclusion and a scalar path
$\varepsilon=t\,v$, $t\ge0$.
\end{proposition}
\begin{reading}
With sign beliefs alone and no regression, the bound is an LP over all unit tables
that fit the margins, and it is exactly right for real-valued tables, conditional
on the observed margins. With whole numbers of ballots it can differ.
\end{reading}

\begin{proposition}[Integer counts, two rows, margins only]\label{prop:integer}\status{proved}
Let $R=C=2$ with no unit-level restrictions, let $N_1$ and $N_0$ be the two
group totals, and $N=N_1+N_0$. The set of gaps $D=B_1-B_2$ attained by integer
tables is the Duncan--Davis interval for $D$ intersected with a lattice of step
$1/N_1+1/N_0=1/(N\mu(1-\mu))$, and both endpoints are attained. Under sign rows
on one column, the integer set is the continuous set of
Proposition~\ref{prop:tworowsign} intersected with that lattice, so the
endpoints differ by less than one step.
\end{proposition}
\begin{reading}
For $N=10^6$ people with $\mu=0.3$ the step is $4.8\times10^{-6}$, so whole
ballots matter only for tiny tables or near-ties. In tiny tables they do matter:
three units of four people with group-1 counts 1, 2, 3 and outcome counts 0, 1, 0
admit the continuous CR set $\{0\}$, but no integer table satisfies CR. With
unit-level restrictions, several columns or $R\ge3$, the integer--continuous gap is
not bounded by this argument; that case is open.
\end{reading}

\begin{proposition}[Two rows]\label{prop:tworowsign}\status{proved}
Let $R=2$, $X=X_1$, $D=B_1-B_2$, $\gamma=\Var_p(X)/(\mu(1-\mu))$, and let
$[D^-,D^+]$ be the Duncan--Davis interval for $D$. Assume $\Var_p(X)>0$, continuous
tables, no unit-level restrictions and no sign rows other than those named. Write
$D_{\ER}$ and $D_{\NM}$ for the one-vs-rest regression slope and neighbourhood
gap.
(a) $\delta_W\ge0\iff D\le D_{\ER}$, and $\delta_B\ge0\iff D\ge D_{\NM}$.
(b) $\{\delta_W\ge0\}$ gives $[D^-,\min\{D^+,D_{\ER}\}]$, empty iff $D_{\ER}<D^-$;
symmetrically, $\{\delta_W\le0\}$ is empty iff $D_{\ER}>D^+$.
(c) $D_{\NM}\in[D^-,D^+]$ always, so a sign on $\delta_B$ is never refuted in this
problem.
(d) $D_{\NM}=\gamma D_{\ER}$ with $\gamma\in(0,1]$. If $\gamma<1$ and
$D_{\ER}\ne0$, exactly one CR branch is non-empty. If $\gamma=1$ (every unit with
positive weight is pure), $D$ is point-identified and both branches are non-empty
with the same gap image $\{D_{\ER}\}$; if $D_{\ER}=0$, both have gap image $\{0\}$.
The CR set is $[D_{\NM},\min\{D^+,D_{\ER}\}]$ if $D_{\ER}\ge0$ and
$[\max\{D^-,D_{\ER}\},D_{\NM}]$ otherwise; under CR,
$\sgn D=\sgn D_{\ER}$ when $D_{\ER}\neq0$.
\end{proposition}
\begin{reading}
In this base problem (two groups, real-valued tables, nothing else imposed) only a
within sign can be contradicted by the margins. A between sign and CR are always
feasible, because ``every group votes like its unit'', $\beta_{i1}=\beta_{i2}=T_i$,
fits them. That non-refutability belongs to this base system only. Under a linear
rung, CR can be rejected jointly with it even when the rung alone fits
(Remark~\ref{rem:LCEexcludesCR}); with integer counts CR can be rejected outright
(Proposition~\ref{prop:integer}); and known cells or unit orderings can remove the
equal-rates witness. In R$\times$C the equal-rates table satisfies every weak
between sign and the CR \emph{union}, but only the branch that matches its within
association.
\end{reading}
\begin{sketch}
By Lemma~\ref{lem:M1}(c) with two rows, $\delta_W=(D_{\ER}-D)\Var_p(X)$ and
$\delta_B=(D-D_{\NM})\mu(1-\mu)$, so each sign is a half-line in $D$, and $D$ ranges
over the whole interval $[D^-,D^+]$. The equal-rates table gives $D=D_{\NM}$. The
identity $D_{\NM}=\gamma D_{\ER}$ with $\gamma\in(0,1]$ shows that the two
same-sign branches of CR overlap only when $\gamma=1$ or $D_{\ER}=0$.
\end{sketch}

\begin{proposition}[Joint restrictions]\label{prop:joint}\status{proved}
Imposing sign restrictions on several columns or rows jointly gives a subset of the
feasible set of any single one of them.
\end{proposition}
\begin{reading}
Joint beliefs are weakly tighter, and often strictly: in $2\times3$ simulations,
signing the within association of two columns jointly narrowed column~1 strictly
in 75 of 139 populations. With two complementary columns, the \emph{same} weak
within sign on both forces both associations to zero and so $D=D_{\ER}$. This is
why the package solves one LP rather than applying $2\times2$ closed forms column
by column.
\end{reading}

\subsection{Signs under a linear rung}

Under a linear rung the model means, for the fixed design, can be written
$m_i=m_{\rm id}(V_i)+L_ik$, where $m_{\rm id}$ is identified, $L_i$ is fixed by the
design, and $k$ stacks the rung's free coordinates: $K$ under M0, the
$K^{(d)}$ under M2, and the common $K$ under M1, with the district shifts in the
identified part.

\begin{proposition}[Sign rows under linear rungs]\label{prop:Krows}\status{proved; (c) for M0 only}
(a) Under every linear rung, $\delta^*_B(r,s)$, $\delta^{r*}_B$, $\Delta^*$ and
$\delta^{r*}_W$ are affine in $k$, so each sign restriction is one linear row in
the LP and CR is a union of two LPs.
(b) If $m_{rc}(x,d)\ge m_{sc}(x,d)$ for every district $d$ and every $x$ in that
district's support, then $\delta^*_B(r,s)\ge0$, pooled over districts.
(c) For $R=2$ under M0, $\delta^*_W=a(w_1+d_1)+b\,w_1$ with $a,b\ge0$, where $w_1$
and $w_1+d_1$ are the two groups' contextual slopes in JKST's notation, so
non-negative slopes imply $\delta^*_W\ge0$. Under the district rungs this
implication fails for the pooled association.
\end{proposition}
\begin{reading}
Under a linear rung, sign beliefs become cheap extra constraints on the unknown
$K$. But a pooled sign under district models is one restriction across all
districts. Positive slopes in every district do not make the pooled within
association positive, because district baselines can fall as composition rises:
by the law of total covariance, the pooled association has a between-district term
with no sign. Example: six equally weighted units with $X=.1,.2,.3$ in district~1
and $.7,.8,.9$ in district~2, both rows' means $.8+.1x$ in district~1 and $.2+.1x$ in
district~2. The slopes are positive, the within-district association is
$+0.00064$, and the pooled one is $-0.04927$.
\end{reading}

\begin{remark}[Realised versus model-implied associations]\label{rem:whichdelta}
N1 uses the realised associations, the linear rows use $\delta^*$. Under (S2),
$\delta-\delta^*=\sum_i\sum_rc_{ir}e_{ir}$ is a design-weighted sum of unit
residuals with mean zero given the design. Under A6 its standard deviation is at
most $\tfrac12\{\sum_i(\sum_r|c_{ir}|)^2\}^{1/2}$; it is small only when no unit
keeps a non-negligible coefficient. A mean condition alone gives no rate.
Intersecting N1 with a linear rung therefore mixes two readings, and the package
reports them separately.
\end{remark}

\subsection{The sign module under the pooled model}\label{sec:msign}

The results of this subsection are proved for the pooled model M0 only. The sign
and CR rows are valid linear rows under the district rungs
(Proposition~\ref{prop:Krows}(a)), but sharpness and inference for them under the
district rungs are open (Section~\ref{sec:open}). The package therefore labels
the module \emph{experimental}.

View $\ThI$ and $\ThO$ in the free coordinates $k$. The target is a model-implied
linear $\varphi(k)=\varphi_0+\omega^\top k$, for example $\Bt_{rc}$ or the gap
$\tilde D=\Bt_1-\Bt_2$. An \emph{admissible restriction} is a set
$C=C(P)\subseteq\R^d$ that is (C1) closed and convex and (C2) depends on $P$ only
through the design and $Q(P)$. Every sign on $\delta^*_B(r,s)$, $\delta^{r*}_B$,
$\Delta^*_r$ or $\delta^{r*}_W$, mean-level ordering on $A$, and the local rows of
Proposition~\ref{prop:sandwich} are admissible and polyhedral. A \emph{union
restriction} is $C_+\cup C_-$ with admissible $C_\pm$; CR for a pair $(g_B,g_W)$
of affine functions of $k$ is $C_\pm=\{\pm g_B\ge0,\pm g_W\ge0\}$.

\begin{theorem}[Sharpness with added restrictions]\label{thm:R3}\status{proved, pooled M0}
Assume A1--A5, $P(X\in A)=1$ and $\affdim(A)=R-1$.
\begin{enumerate}[label=(\alph*)]
\item For admissible $C$, the sharp identified set for $\varphi$ under LCE and $C$
is $[\min_{\ThI\cap C}\varphi,\max_{\ThI\cap C}\varphi]$, or empty.
\item $\ThO\cap C=\bigcup_{P'\in\mathcal P(P)}(\ThI(P')\cap C(P'))$. So the
interval $\varphi(\ThO\cap C)$ is sharp for the regression and the design. It is a
finite LP when $C$ is polyhedral, it equals (a) when $\ThI=\ThO$, and it is weakly
outer otherwise.
\item For a union, the identified set is $\varphi(\Theta\cap C_+)\cup\varphi(\Theta\cap C_-)$
for $\Theta\in\{\ThI,\ThO\}$: at most two intervals. The hull is sharp iff they
intersect, and a common point of $\Theta\cap C_+\cap C_-$ certifies that they do
(sufficient only).
\item For one-vs-rest CR, with $\Var_p(X_r)>0$,
$\gamma_r=\Var_p(X_r)/(\mu_r(1-\mu_r))$, $\ER^{r*}=\Cov_p(X_r,\tau(X))/\Var_p(X_r)$ and
$D^{r*}=\Bt_r-\Bt_{\neg r}$,
\[
C_+\cup C_-=\{k:\ D^{r*}(k)\in\hull\{\gamma_r\ER^{r*},\ \ER^{r*}\}\},
\]
which is two linear rows. For $R=2$ this is the only CR, so in $2\times2$ CR is a
single LP, and under it $\sgn\tilde D=\sgn\ER^*$.
\item Refutation is joint: if $P$ is generated by a model in LCE $\cap\,C$, the
true $k_0\in\ThI\cap C\subseteq\ThO\cap C$; so an empty population LP rejects the
conjunction, not either part alone.
\end{enumerate}
\end{theorem}
\begin{reading}
Adding a sign or ordering belief to the pooled linear rung keeps the bound as
tight as the regression allows: nothing narrower can be learned from the
regression and the design. The full law of the data could rule out more. For the
default one-vs-rest comparison, reinforcement becomes a single band in the gap,
so no union is needed. If the bound is empty, it is the combination that is
contradicted, not the belief on its own. A belief that also has a model-free
meaning can sometimes be tested on its own, and conversely a linear rung can
exclude a belief that is never refutable model-free.
\end{reading}
\begin{sketch}
(a) The identified set is the image of the compact convex set $\ThI\cap C$ under an
affine map; (C2) makes $C$ the same for every law consistent with the data.
(b) Apply the construction of Proposition~\ref{prop:sharpreg} to each point of
$\ThO\cap C$. (d) Apply Lemma~\ref{lem:M1}(c) to the model means:
$\delta^{r*}_W=\Var_p(X_r)(\ER^{r*}-D^{r*})$ and
$\delta^{r*}_B=\mu_r(1-\mu_r)(D^{r*}-\gamma_r\ER^{r*})$; the union of the two
same-sign cases is the hull.
\end{sketch}

\begin{remark}[Full law versus regression]\label{rem:fullreg}
Let $X\in\{1/4,1/2,3/4\}$ with equal probability and equal sizes, and $T$ a fair
binary total at every composition. The full law forces both rows' rates to equal
$T$ in every unit, so $k=0$. The regression alone admits $m_1=.5+k(1-x)$,
$m_0=.5-kx$. For $C=\{k\in[.2,.4]\}$ the full-law set of (a) is empty, while the
regression-only gap set of (b) is $[1/6,1/3]$.
\end{remark}

\begin{remark}[A linear rung can exclude CR although it fits]\label{rem:LCEexcludesCR}
Let $X\in\{0,1/4,1\}$ with equal probability, $A=[0,1]$, $m_0=1-x$,
$m_1=(49/55)x$ and deterministic $T$. The vertex rows force $k=0$, so LCE fits at a
single point, where $\ER=0$ and $\tilde D=-261/1925$. CR requires
$\tilde D\in\{0\}$, so both branches are empty, although LCE alone is not
rejected. In the model-free continuous problem CR is still feasible. Any statement
that a belief ``cannot be refuted'' must name its base system.
\end{remark}

For CR pairings other than W1, part (d) fails: the two sign conditions are on
different affine functions of $k$, so both branches can be non-empty with
disjoint images. An adversarial search found gaps of up to about 31\% of the hull
width. For the pairwise version only $\delta^*_B(r,s)$ is pairwise; $\Delta^*_r$
pools all groups, including those outside the pair.

\begin{proposition}[When does a restriction tighten the bound?]\label{prop:R2}\status{proved}
Let $F$ be the current non-empty compact feasible set, $U=\max_F\varphi$,
$L=\min_F\varphi$, $F^*_U$ and $F^*_L$ the optimal faces, and $N$ an admissible
restriction with $F\cap N\neq\emptyset$ and new endpoints $U_N,L_N$.
\begin{enumerate}[label=(\alph*)]
\item $U_N<U$ iff $F^*_U\cap N=\emptyset$, and $L_N>L$ iff $F^*_L\cap N=\emptyset$.
For polyhedral $N$ each is one feasibility LP, which must carry the target
intercept: $\{Gk\le g,\ \omega^\top k\ge U-\varphi_0,\ k\in N\}$.
\item For one row $a^\top k\le b$: $U_N<U$ iff $\min_{F^*_U}a^\top k>b$, iff every
optimal dual multiplier of the row is positive. A positive multiplier at one
solver optimum does not certify tightening.
\item $N$ is implied ($F\subseteq N$) iff it tightens no linear target.
\item In $2\times2$ under LCE with $\Var_p(X)>0$ and $\gamma<1$, each of the four sign
rows is either implied or tightens every target that is non-constant in the
unknown slope; the closed forms are those of Proposition~\ref{prop:G2}.
\item For a stack of rows, a row's contribution depends on what is already
imposed. Zero gain when added last means zero incremental gain for this target,
not that the row is implied.
\end{enumerate}
\end{proposition}
\begin{reading}
A belief narrows a quantity only if it removes every table that attains that end
of the bound. The package tests this with one feasibility LP per end and never
trusts solver dual values: a row duplicated at an optimal vertex can carry a
positive multiplier without moving anything. A belief can narrow the set of
possible tables without narrowing the quantity of interest; with three groups this
happened in 28\% of simulated cases, against 8\% implied and 64\% tightening.
\end{reading}
\begin{sketch}
Adding a restriction can only lower the maximum. If some maximiser survives, the
maximum is unchanged; if none does, the restricted problem is compact and
non-empty and attains a strictly smaller value. For one row, concavity of the LP
perturbation function ties this to every optimal multiplier being positive. A
hyperplane separating a point of $F\setminus N$ from $N$ gives a target that
tightens, which proves (c).
\end{sketch}

\begin{proposition}[JKST plus signs in $2\times2$]\label{prop:G2}\status{proved}
Under LCE in $2\times2$ with $\Var_p(X)>0$ and $\gamma<1$, let $[L_J,U_J]$ be the
JKST interval for the gap $\tilde D$. Then $\delta^*_W\ge0$ gives
$[L_J,\min(U_J,\ER^*)]$; $\delta^*_W\le0$ gives $[\max(L_J,\ER^*),U_J]$;
$\delta^*_B\ge0$ gives $[\max(L_J,\NM^*),U_J]$; $\delta^*_B\le0$ gives
$[L_J,\min(U_J,\NM^*)]$; and CR gives $[L_J,U_J]\cap\hull\{\NM^*,\ER^*\}$.
Reversed endpoints mean the set is empty.
\end{proposition}
\begin{reading}
Adding a sign belief to JKST cuts its interval at one of two data-computed
numbers. In 2000 random LCE designs with true signs, the within sign tightened the
JKST interval in 83\% of designs (median width ratio 0.70), the between sign in
89\% (0.67), and CR in 80\% of the designs where CR holds (0.69). A sign is
redundant exactly when the rung alone already implies it, that is, when the
interval for that association implied by the rung lies on the believed side of
zero. Under LCE, in $2\times2$ the pair
$(\delta_W,\delta_B)$ lies on a line segment, so re-imposing the rung's own
implications on the associations adds nothing.
\end{reading}

\paragraph{Orderings and local beliefs.} For one column and any $R$, write
$q_{rs}(x)=(Qx)_r-(Qx)_s$, which is identified, and $g_{rs}(x;k)=(Kx)_r-(Kx)_s$,
linear in $k$. The between-group gap at $x$ is $O_{rs}(x)=(Hx)_r-(Hx)_s$, and the
composition-weighted contextual derivative when a member of group $s$ is replaced
by one of group $r$ is $L_{rs}(x)=(x^\top H)_r-(x^\top H)_s$.

\begin{proposition}[Sandwich]\label{prop:sandwich}\status{proved}
(a) $O_{rs}(x)+L_{rs}(x)=2q_{rs}(x)=\partial_{e_r-e_s}\tau(x)$. (b)
$O_{rs}=q_{rs}+g_{rs}$ and $L_{rs}=q_{rs}-g_{rs}$. (c) $O_{rs}\ge0$ and
$L_{rs}\ge0$ at $x$ iff $|g_{rs}(x;k)|\le q_{rs}(x)$, which is possible only if
$q_{rs}(x)\ge0$, a data test.
\end{proposition}
\begin{reading}
An ordering belief (``group $r$ is at least as supportive as group $s$'') and a
local contextual belief restrict the same unknown from opposite sides, so together
they confine it to an interval whose half-width is an identified slope.
\end{reading}

\begin{proposition}[Which directions a pairwise belief can reach]\label{prop:rank}\status{proved}
(a) One pair $(r,s)$ restricts an $(R-1)$-dimensional subspace of the column's
$K$. (b) Pairs forming a star $\{(1,s):s\ne1\}$ reach all $R(R-1)/2$ directions.
(c) Ordering and local rows (Proposition~\ref{prop:sandwich}) on a star bound the
column's $K$ without Duncan--Davis rows.
\end{proposition}
\begin{reading}
With three or more groups, each pairwise comparison can restrict only $R-1$ of
the unknown directions. Enough well-chosen pairs reach all of them; whether any
is actually cut depends on the beliefs and the data. This is the precise sense in
which more groups give more room for beliefs to bind. That every spanning tree of
pairs, not only a star, reaches all directions has been checked numerically for
chains with $R=3,4,5$ but is not proved, and what happens when the support of
compositions is lower-dimensional is open.
\end{reading}

\begin{proposition}[Nesting of beliefs in $2\times2$ under M0]\label{prop:PN}\status{(a), (b) proved}
Let $a=\mu\Var_{w^1}(X)$, $b=(1-\mu)\Var_{w^2}(X)$ and $x^*=a/(a+b)$. (a) Signed
slopes ($w_1\ge0$, $w_1+d_1\ge0$) imply the weighted-local belief $w_1+d_1x\ge0$ on
the hull. (b) $\delta^*_W=(a+b)(w_1+d_1x^*)$, so the weighted-local belief implies
$\delta^*_W\ge0$ whenever $x^*$ lies in the composition hull.
\end{proposition}

\begin{conjecture}[The balance point lies in the hull]\label{conj:C1}
In $2\times2$, $x^*$ always lies in the composition hull.
\end{conjecture}
\begin{reading}
In the pooled $2\times2$ model the weighted-local belief (``replacing a member of
one group by one of the other does not lower the composition-weighted rates, at
every composition'') sits between signed slopes and a sign on the within
association. The second step needs $x^*$ to lie in the observed range.
Conjecture~\ref{conj:C1} says it always does: there was no violation in 2000
random designs and an adversarial search over 50{,}000 discrete designs, but there
is no proof. The chain fails for pooled M1 (Proposition~\ref{prop:Krows}).
\end{reading}

\subsection{Monotone nonparametric contextual effects (N2)}\label{sec:N2}

N2 drops linearity and keeps only direction. It is defined for two rows on the
finite empirical design law: the support of $X$ (the row-1 share) is the set of
distinct observed compositions $x_1<\dots<x_{J_X}$, each with positive mass. Model
N2($\sigma$) assumes (S1), $\E[\beta\mid V]=m(X)$ for a Borel
$m:[0,1]\to\Bset$, and, for designated cells $(r,c)$ with directions
$\sigma_{rc}\in\{+1,-1\}$, that $\sigma_{rc}m_{rc}$ is non-decreasing on the
support. The regression $\tau(x)=\E[T\mid X=x]$ is identified at each support
point.

\begin{lemma}[Means lie in the mean polytope]\label{lem:N21}\status{proved}
$m(x)\in S(x,\tau(x))$ at every support point.
\end{lemma}

Let $\mathcal M_O$ be the set of support-point means with $m(x_j)\in S(x_j,\tau(x_j))$
for all $j$ and
$\sigma_{rc}(m_{rc}(x_{j+1})-m_{rc}(x_j))\ge0$ for all adjacent pairs and designated
cells. Bounds on $\Bt_{rc}=\sum_iw_{ir}m_{rc}(X_i)$ are an LP over $\mathcal M_O$.

\begin{proposition}[The N2 outer set]\label{prop:N2}\status{proved, finite support}
(a) Every $m$ allowed by N2($\sigma$) restricts to a point of $\mathcal M_O$, so the
LP bounds are an outer set for $\Bt$; they are not claimed sharp.
(b) Every point of $\mathcal M_O$ extends to a monotone Borel $m$ on $[0,1]$, so on
the complete finite support nothing is lost by restricting to support points. On a
coarser grid this fails, and the grid LP is only a relaxation.
(c) If $\tau$ is quadratic, M0 with signed slopes lies inside N2, for the outer sets.
(d) If $m_{1c}$ and $m_{2c}$ are both non-decreasing, $\delta^*_W\ge0$ for column $c$.
(e) For $R\ge3$ a direction ``in $x_1$'' needs a chart saying which share adjusts,
and (d) has no general analogue.
\end{proposition}
\begin{reading}
N2 assumes only ``no reversals'': each designated group's mean rate moves in one
direction as its share grows, in any shape. It is weaker than straight lines with
signed slopes, so its bound contains theirs, and the two-row shape restriction is
stronger than a pooled covariance sign, which permits reversals. If a designated
mean reverses direction somewhere (a U-shape or a tipping point), the order rows
exclude the true mean function, and the bound can be narrow, non-empty and wrong;
an empty set is the only warning, and plotting the fitted regression for
reversals is a partial check. N2 is not in the package.
\end{reading}

\begin{conjecture}[The sharp N2 set]\label{conj:N2}
The sharp N2 set is the set of $m$ with $m(X)\in\E[S(X,T)\mid V]$ $P_V$-almost
surely, plus the shape restrictions, where the expectation is the conditional
selection expectation given $V=(X,N)$.
\end{conjecture}
\begin{reading}
When unit size carries information, conditioning on composition alone admits mean
rates that the full model excludes. Fix $X=1/2$ and two equally likely sizes; at
the first size $T$ is 0 or 1 with equal probability, at the second $T=1/2$ surely.
At the first size both rows' rates must equal $T$, so the model forces
$m(1/2)=(.5,.5)$; the polytope LP, which conditions on $X$ only, admits
$(.75,.25)$. Under the selection-expectation reading the characterisation is
sufficient on a finite support: for each design atom take a selection whose
conditional mean is the common candidate and glue the selections across atoms.
Necessity in general is not written, so the statement stays a conjecture. A
procedure that uses only $\E[T\mid X]$ does not attain the set when size is
informative.
\end{reading}

\begin{proposition}[Local identity]\label{prop:N23}\status{proved, as a calculus identity}
Let $m_{1c},m_{2c}$ be differentiable at $x$, and let
$\tau_{m,c}(z)=z\,m_{1c}(z)+(1-z)m_{2c}(z)$ be the structural aggregate,
$D_c(x)=m_{1c}(x)-m_{2c}(x)$ and $\delta_{W,c}(x)=x\,m'_{1c}(x)+(1-x)m'_{2c}(x)$.
(a) $\tau'_{m,c}(x)=D_c(x)+\delta_{W,c}(x)$. (b) $\delta_{W,c}(x)\ge0\iff
D_c(x)\le\tau'_{m,c}(x)$. (c) Local CR, $D_c(x)\delta_{W,c}(x)\ge0$, holds iff
$D_c(x)$ lies between 0 and $\tau'_{m,c}(x)$. (d) If both rows are non-decreasing on
a neighbourhood of $x$, the local sign at $x$ is non-negative; monotonicity on the
finite support is not enough, and the local sign everywhere neither implies nor is
implied by the sign of $\delta^*_W$.
\end{proposition}
\begin{reading}
EGGHM's local results are statements about smooth functions, but on a finite
support the data never identify a derivative: adding
$\prod_j(x-x_j)$ to both rows leaves the law unchanged and moves the derivative at
every node. So local sign beliefs cannot be imposed from aggregates without an
identified regression family, and they are not an observable rung here.
\end{reading}

\begin{lemma}[Person-weighted bins]\label{lem:bins}\status{proved}
Let $b=1,\dots,K$ be ordered non-overlapping bins of compositions with
person-weighted masses $P(b)>0$ and within-bin expectations $\E_b$. Put
$\bar x_b=\E_bX$, $\bar m_{1b}=\E_b[Xm_1(X)]/\bar x_b$,
$\bar m_{2b}=\E_b[(1-X)m_2(X)]/(1-\bar x_b)$ and $\bar\tau_b=\E_b[\tau(X)]$, the
model-implied bin mean. Under N2($\sigma$): (i) $\bar m_b\in S(\bar x_b,\bar\tau_b)$;
(ii) the targets are linear in the bin rates with known weights; (iii) monotonicity
is preserved across the ordered bins. The binned LP is an outer relaxation of the
N2 bounds.
\end{lemma}
\begin{reading}
With many compositions, group them into bins, but average rates by weighting the
people in each group, not the units: the ordinary bin average breaks the
accounting identity (it gives 0.25 against the correct 0.3125 in a two-point
example). In samples the observed person-weighted bin mean of $T$ estimates
$\bar\tau_b$ with error, which inference must account for.
\end{reading}

\paragraph{What inference for N2 needs.} The regression at the support points
must be estimated. A weighted polynomial of \emph{prespecified} degree $q$ makes
the constraints linear in its coefficients, and the confidence intervals of
Section~\ref{sec:inference} then apply conditional on the degree. The degree must
not be chosen by cross-validation and then treated as fixed: in a balanced
two-level design with true slope zero, choosing between an intercept-only and a
linear fit by leave-one-out error and then using the nominal 95\% slope interval
gives conditional coverage of about 0.68.
